problem:  
Let \(M^n\) be a closed, simply connected Riemannian manifold with **nonnegative sectional curvature**, equipped with an **isometric fixed point homogeneous circle action**.  
Let \(F \subset \mathrm{Fix}(S^1,M)\) be the maximal-dimensional fixed component, and let \(Q = M/S^1\) denote the Alexandrov orbit space (curvature \(\ge 0\), boundary \(\partial Q = F\)). Define the distance function \(r(x) = \mathrm{dist}_Q(x, F)\) and let \(E \subset Q\) be the **extremal set** where \(r\) attains its maximal value.  

**Problem:**  
It is known that for positively curved \(M\), \(Q\) is an Alexandrov join \(Q \cong F * E\). For nonnegative curvature, extensive structure theory ensures only concavity and extremality properties.  

The open question is:  
> **Does nonnegative curvature together with fixed point homogeneity imply that \(Q\) contains an open dense subset isometric (in the Alexandrov sense) to a join \(F^\circ * E^\circ\), where \(F^\circ \subset F\) and \(E^\circ \subset E\) are open dense smooth strata?**  
> Equivalently, is there a *local-to-global join rigidity phenomenon* for nonnegatively curved fixed point homogeneous orbit spaces—namely, does the infinitesimal join structure near \(F\) (guaranteed by the fixed point homogeneity condition) extend to a global metric join structure on a dense subset of \(Q\)?  

If not, can one construct or preclude examples where the local join model fails to extend, yielding a nonnegatively curved \(M\) whose orbit space \(Q\) exhibits genuinely new Alexandrov geometry local-to-global obstruction?  

Why is it a "good" problem:  
This problem probes the *precise frontier of rigidity phenomena* in the transition from positive to nonnegative curvature under group symmetry.  

- **Profound structural insight:** Positively curved fixed point homogeneous manifolds exhibit global Alexandrov join structure in their orbit spaces—this is a cornerstone fact in classification results. The proposed question asks whether the same principle survives under weaker curvature, which lies at the heart of understanding when **local geometric symmetry forces global metric rigidity**.  
- **Bridging two domains:** It unites **Riemannian group action theory** (fixed point homogeneity and orbit decomposition) with **Alexandrov geometry** (extremal subsets, local join structures, convexity and distance concavity). Resolving this would likely require blending smooth and synthetic methods, creating a conceptual bridge between global smooth geometry and metric curvature theory.  
- **Simple but deep:** The statement involves only natural geometric quantities (\(F\), \(E\), distance function, join notion), yet it encodes the delicate **local-to-global propagation of curvature symmetry**, a phenomenon known to hold in positive but not clearly understood in nonnegative curvature.  
- **Transformative potential:** A positive answer would establish new global rigidity for nonnegative Alexandrov quotient spaces—essentially extending the classical **Grove–Searle join rigidity** beyond positive curvature. A negative answer would reveal new geometric flexibility, identifying how *collapse* or *flat directions* destroy join propagation, reshaping how curvature and symmetry interact in Riemannian geometry.

Thus, it qualifies as a “good” problem: it is compactly and cleanly stated, centered on a pivotal rigidity question, and positioned to advance the conceptual boundary between Alexandrov and Riemannian geometry in the study of symmetry and curvature.